\subsection{Edge Insertion}
\label{sec:edgeInsertion}

As aforementioned, inserting an edge may trigger the updates of several CL-tree nodes as well as their structures.
We illustrate this by Example~\ref{eg:edgeInsertion}.

\begin{example}
\label{eg:edgeInsertion}
Consider the graph in Figure~\ref{fig:advancedIndex}. If we insert an edge ($H$, $G$) as shown in Figure~\ref{fig:coreNumber}, the core number of vertex $H$ increases to 2 and we need to move it down to a node in the lower level. If we insert an edge ($G$, $I$), the connectivity of some vertices changes as shown in Figure~\ref{fig:ConnectivityExmp} and thus the corresponding subtrees are merged as a new one.
\end{example}

\begin{figure}[ht]
    \centering
    \mbox{
        \subfigure[core number]{
            \includegraphics[width=.4\columnwidth]{figures/movedownExmp}
            \label{fig:coreNumber}
        }
        \hspace{2ex}
        \subfigure[connectivity]{
            \includegraphics[width=.4\columnwidth]{figures/connectiveExmp}
            \label{fig:ConnectivityExmp}
        }
    }
    \caption{The core number and connectivity change.}
    \label{fig:connectivity}
\end{figure}

To maintain the CL-tree for inserting an edge, we propose an algorithm called {\tt insertEdge}.
The main idea is that, we first find vertices whose core numbers change, then change their positions in the CL-tree, and merge some subtrees. Let $V^+$ be the set of vertices whose core numbers increase after inserting an edge ($u,v$).
We summarize the main steps of {\tt insertEdge} as follows.

$\bullet$ \textbf{Step 1:} Compute $V^+$;

$\bullet$ \textbf{Step 2:} Move down vertices of $V^+$;

$\bullet$ \textbf{Step 3:} Merge subtrees.

We now elaborate these steps one by one.

\textbf{Step 1: Compute $V^+$.}
Inserting an edge only affects the core numbers of a small number of vertices~\cite{kcoreUpdate,sariyuce2016incremental}. We first give a definition, a theorem and a lemma proposed in this paper.

\begin{definition}[\cite{kcoreUpdate}]
\label{df:inducedgraph}
Given a graph $G$ and a vertex $v$, the induced core subgraph of $v$, denoted as $G_v$, is a connected subgraph containing $v$ and the core numbers of all vertices in $G_v$ equal to $core_{G}[v]$.
\end{definition}

Notice that, the sets of vertices in $G_u$ ($G_v$) are actually subsets of vertices in $p_u.vertexSet$ ($p_v.vertexSet$),
where $p_u$, $p_v$ denote the nodes that contain $u$, $v$.

\begin{theorem}[k-core update theorem\cite{kcoreUpdate}]
\label{thrm:kcoreupdate}
Given a graph $G$ and two vertices $u$ and $v$. After inserting or deleting an edge $(u$,$v)$ in $G$, we have that,

$\bullet$ If $core_G[u] > core_G[v]$, only the core numbers of vertices in $G_v$ may need to be updated.

$\bullet$ If $core_G[u] < core_G[v]$, only the core numbers of vertices in $G_u$ may need to be updated.

$\bullet$ If $core_G[u] = core_G[v]$, only the core numbers of vertices in the union of $G_u$ and $G_v$, \textit{i.e.,} $G_{u\cup v}$ may need to be updated.
\end{theorem}

\begin{lemma}[\cite{kcoreUpdate}]
\label{lm:kcorelemma}
After inserting (deleting) an edge, the core number of any vertex in $G$ increases (decreases) by at most 1.
\end{lemma}

By above theorem and lemma, we can conclude that only a small number of vertices need to change their core numbers.
In specific, we can first find node $p_u$ ($p_v$) and then compute the vertex set $V^+$ in which vertices's core numbers increase by 1 using the algorithm in~\cite{kcoreUpdate}.

\textbf{Step 2: Move down vertices of $V^+$.}
Let $p$ be the node containing $V^+$ and $c$=min\{$core_G[u], core_G[v]$\}). Since the core numbers of vertices in $V^+$ increase by 1 (from $c$ to $c$+1), we need move them down to nodes in the lower level.
During the moving down process, we may also need to reorganize $p$'s child nodes. Let us illustrate this by Example~\ref{eg:goDown}.

\begin{figure}[ht]
    \centering
    \includegraphics[width=1.04\linewidth]{figures/movedownEmp}
    \caption{An example of the tree index update.}
    \label{fig:movedownEmp}
\end{figure}


\begin{example}
\label{eg:goDown}
Consider a graph in Figure~\ref{fig:movedownEmp}(a) and its CL-tree in Figure~\ref{fig:movedownEmp}(b).
Let us insert an new edge (8, 11). We first get $V^+$=\{8, 11, 23\} and $c$=2.
Next, we move them down from $r_1$ to $r_3$.
Besides, we have to merge $r_2$ into $r_3$ and place $r_4$ as $r_3$'s child node,
since their connectivity changes after the insertion.
The updated CL-tree is depicted in Figure~\ref{fig:movedownEmp}(c).
%Consider that an edge ($8,11$) is inserted in Figure~\ref{fig:movedownEmp}(a). Before insertion, the CL-tree index is shown in Figure~\ref{fig:movedownEmp}(b). After insertion, vertices $8$, $11$ and $23$ form $V^+$ and need to increase the core number from 2 to 3. Thus they need to move down a level and the tree index changes accordingly. Nodes $r_2$, $r_3$ whose core number are 3 are traced by the neighbors of the vertices ($11$ and $23$) and therefore $r_2$, $r_3$ are merged as a new node (see the arrows in Figure~\ref{fig:movedownEmp}(b)). Another node $r_4$ which is connected through $8$ is also affected and re-linked as a child node of the new node $r_2$ because the core number of $r_4$ is 4. Other nodes remain unchanged. Finally the updated tree index is presented in Figure~\ref{fig:movedownEmp}(c).
\end{example}

%We divide $p$'s child nodes into two sets $Z_1$ and $Z_2$, where nodes in $Z_1$ have a core number of $c$+1 and nodes in $Z_2$ have a core number of $c$+2 or more.

Clearly, moving down vertices of $V^+$ from $p$ to $p$'s child node (denoted by $p'$) may change the connectivity of $p$'s child nodes.
Consider a specific vertex $a$$\in$$V^+$ and we initialize two empty sets $B_1$ and $B_2$.
For each of $a$'s neighbor $b$ whose $core_G[b]$$\textgreater$$c$, we first find the node $p_b$ containing $b$,
and then trace it up from $p_b$ along the CL-tree until a child node of $p$, denoted by $o_b$.
If $o_b$ has a core number of $c$+1, we put it into $B_1$;
Otherwise, we put it into $B_2$.
Then, after moving down vertices of $V^+$,
nodes in $B_1$ should be merged into $p'$ and nodes in $B_2$ will be child nodes of $p'$.


%Some child nodes of $p$ which are connected through vertices in $V^+$ may also need to be updated. These child nodes can be divided into two types. (1) The core number of the node is $c$+1; (2) The core number of the node is larger than $c$+1. If $c=core_G[v]$, for example, after increasing the core number of the vertex $v$ to $c$+1, all ($c$+1)-$\widehat{core}$ which are connected with $v$ should be merged as a new ($c$+1)-$\widehat{core}$. Reflected in the tree index, the first type of child nodes mentioned above will be merged as a new tree node. In addtion, if there exists the second type of child node, this node will become the new child node of the new node whose core number is $c$+1. This process should be done for each vertex in $V^+$. We give the Algorithm~\ref{alg:moveDown} and illustrate this in Example~\ref{eg:goDown} afterwards.

\begin{algorithm}
\caption{move down vertices: {\tt moveDown}}
\label{alg:moveDown}
\footnotesize{
\algrenewcommand{\algorithmiccomment}[1]{\hskip3em$//$ #1}
\begin{algorithmic}[1]
    \Function{moveDown($V^+$, $p$)}{}
    \If{$V^+$=$\emptyset$}
        \Return $p$;
    \EndIf
    \State $P \gets \emptyset$;
    \State update $p$ using $V^+$;
    \For {each $a \in V^+$}
        \For{each $b\in a$'s neighbor vertices}
            \If{$core_G[b]>c$ and $b\notin V^+$}
                \State locate node $p_b$;
                \State run \Call{TRACE($p_b$)}{}, and update $P$;
            \EndIf
        \EndFor
    \EndFor
    \State $p_{max}\gets$ a node of $P$, which has a core number of $c$+1
        and its $vertexSet$ is the largest among all nodes of $P$;

    \If{$p_{max} = $ null }
        \State create a new node $p'$;
        \State update $p'$;
        \State add $P$ to $childList$ of $p'$;
    \Else{}
    \State add $V^+$ to $p_{max}.vertexSet$;
         \For{each $p_i \in P$}
            \If{$p_i.coreNum$ = $c+1$}
                \State merge $p_i$ to $p_{max}$;
            \Else{}
                \State  add $p_i$ to $childList$ of $p_{max}$;
            \EndIf
         \EndFor
         \State $p' \gets p_{max}$;

    \EndIf
    \State update vertex-node map;
     \If{$p.vertexSet=\emptyset$}
            \State add $ \{p.childList-P\}$ to $childList$ of $p.father$;
         \EndIf
    \State \Return $p'$;
    \EndFunction
\end{algorithmic}}
\end{algorithm}

Algorithm~\ref{alg:moveDown} presents {\tt moveDown}.
If $V^+$$\neq\emptyset$, we first initialize a node set $P$ (line 3).
Then, we remove $V^+$ from $p.vertexSet$ and update $p.invertedList$ (line 4).
$\forall$$a\in V^+$, we enumerate $a$'s neighbor $b$ whose $core_G[b]$$\textgreater$$c$,
locate $p_b$, trace up from $p_b$ to find $p_b$'s ancestor node $o_b$ which is a child node of $p$,
and put $o_b$ into $P$ (lines 5-9).
Let the node which has the largest size with core number being $c$+1 in $P$ be $p_{max}$ (line 10).
Next, if $p_{max}$=$null$, we need to create a new child node of $p$ (lines 11-14);
otherwise, we merge and reorganize $p$'s child nodes (lines 15-22).
Finally we return node $p'$ (line 26), which will be used later.

%In Algorithm~\ref{alg:moveDown}, if $V^+$ is not an empty set, we first initialize a node set $P$ and a node $p_{max}$ (line 3). After computing $V^+$, we need to remove $V^+$ from $vertexSet$ of $p$ and update $invertedList$ of $p$ (line 4). For each vertex $a\in V^+$, we locate $a$'s neighbor vertex $b$ whose core number is larger than c in the tree at the node $p_b$, then we trace up from $p_b$ with $fatherNode$ to find $p_b$'s ancestor node which is also the child node of $p$ and add it to $P$. Simultaneously, we mark the child node whose core number is $c$+1 and contains the most vertices as $p_{max}$ (lines 5-9). Next, if $p_{max}$ does not exist, a new node $p'$ will be created (line 11). Vertices in $V^+$ will be added to $vertexSet$ of $p'$ and all child nodes in $P$ will also be linked to $p'$ (lines 12-13). If $p_{max}$ exists, vertices in $V^+$ will be added into $vertexSet$ of $p_{max}$. All the first type of child nodes in $P$ will be merged to $p_{max}$, and the second type of child nodes in $P$ will be linked to $p_{max}$ (line 15-20). Then we update the vertex-node map (line 22). Note that after removing $V^+$, if $p.vertexSet$ is empty, we update the child list of $p.father$ (lines 23-24). Finally the updated node will be returned (line 25).
%Note that in the process of merging nodes, we merge nodes to the one that contains the most vertices because it is more efficient.



\textbf{Step 3: Merge subtrees.}
Recall in Figure~\ref{fig:ConnectivityExmp}, after inserting ($G$, $I$), the corresponding subtrees, which correspond to the $k$-$\widehat{core}$s containing $G$ and $I$ are merged into one subtree. The process of merging subtrees starts from the tree nodes which contain $G$ and $I$, and ends at their common ancestor node.
The merging process is guaranteed by Lemmas~\ref{lemma:mergetree} and \ref{lemma:relevantNode}.

\begin{lemma}
\label{lemma:mergetree}
After inserting an edge between two vertices, the maximum numbers of disconnected $k$-$\widehat{core}$s
which need to be merged is two.
\end{lemma}
\begin{proof}
It can be proved directly by contradiction.
\end{proof}

%\begin{figure}[ht]
%    \centering
%    \includegraphics[width=0.4\linewidth]{figures/MGlemma}
%    \caption{Three separate $k-\widehat{core}$.}
%    \label{fig:MGlemma}
%\end{figure}
%
%\begin{proof}
%\label{prf:proof}
%Since
% One subtree whose core number of the root is $k$, for instance, represents a $k$-$\widehat{core}$. Lemma~\ref{lemma:mergetree} means, in other words, after inserting an edge between two vertices, there are at most 2 $k$-$\widehat{core}$ that need to be merged.
%
% Figure~\ref{fig:MGlemma} presents that $G_1$, $G_2$, $G_3$ are three $k$-$\widehat{core}$, and $u \in G_1$, $v \in G_2$, $w \in G_3$ . Suppose that, after inserting an edge($u$,$v$), $G_1$ and $G_2$ is connected and need to be merged. From the definition of k$-\widehat{core}$, every vertex in $G_1$ can reach every vertex in $G_2$ along the path ($u$, $...$, $v$). If $G_3$ is also affected by the insertion which means $G_3$ needs to be merged with $G_1$ and $G_2$. There must exist one path ($w$, ..., $u$, ..., $v$). Suppose that the path is between $w$ and $v$, there are two possible paths for $w$ to reach $v$. One is directly through the path ($w$, ..., $v$), the other is through ($w$, ..., $u$) first, and then through ($u$, ..., $v$). Case one means there exists an edge between $G_2$ and $G_3$, which is contradictory to the fact that $G_2$ and $G_3$ are two separate subgraphs. Case two means there already exists one edge before the insertion, and that is also contradictory to the fact.
%
% Therefore, the number of subtrees which need to be merged is less than 3. This completes the proof.
%\end{proof}

\begin{lemma}
\label{lemma:relevantNode}
In the process of merging subtrees, the maximum number of nodes which need to be merged in each level is two.
\end{lemma}
\begin{proof}
It can be proved directly by contradiction.
\end{proof}
%\begin{proof}
%The hypothesis that there exists a third tree node which need to be merged in each level can be proved to be false in the similar way to the proof of Lemma~\ref{prf:proof}. Thus, we do not repeat it.
%\end{proof}

By Lemmas~\ref{lemma:mergetree} and \ref{lemma:relevantNode}, we conclude that,
to merge the subtrees, we can first trace two paths starting from $p_u$ and $p_v$ until their common ancestor in the CL-tree,
and then merge the pairs of nodes on the paths, if their core numbers are the same.

Algorithm~\ref{alg:insertEdge} presents the overall steps of {\tt insertEdge}.
Following Theorem~\ref{thrm:kcoreupdate}, we first compute $V^+$,
and invoke {\tt moveDown} to update these nodes in CL-tree (lines 2-16).
Next, if $p_u'$ and $p_v'$ belong to two disconnected $k$-$\widehat{core}$s, we need to merge the subtrees (lines 17-19).
In detail, we first trace two paths starting from $p_u'$ and $p_v'$ up until one common ancestor.
Then, for each pair of nodes on the paths, if their core numbers are equal,
we merge them as a single node.
Finally, the tree index is updated.
Note that during the above process, the elements of nodes and vertex-node map are also updated.
%the efficient way is always merging the node with few vertices to the one with more vertices.

\begin{algorithm}[h]
\caption{index update algorithm: {\tt insertEdge}}
\label{alg:insertEdge}
\footnotesize{
\algrenewcommand{\algorithmiccomment}[1]{\hskip3em$//$ #1}
\begin{algorithmic}[1]
    \Function{insertEdge($p_u$,$p_v$)}{}
    \If{$p_u.coreNum$=$p_v.coreNum$}
        \State compute  $V_1^+$ in $p_u.vertexSet$;
        \State $p_u' \gets $ \Call{movedown($V_1^+$,$p_u$)}{};
        \State $p_v' \gets p_v$;
        \If {$p_u \neq p_v$}
            \State compute $V_2^+$ in $p_v.vertexSet$;
            \State $p_v' \gets $ \Call{movedown($V_2^+$,$p_v$)}{};
        \EndIf
    \ElsIf{$p_u.coreNum < p_v.coreNum$}
        \State compute $V^+$ in $p_u.vertexSet$;
        \State $p_u' \gets $ \Call{movedown($V^+$,$p_u$)}{};
        \State $p_v' \gets p_v$;
    \Else
        \State compute $V^+$ in $p_v.vertexSet$;
        \State $p_v' \gets $ \Call{movedown($V^+$,$p_v$)}{};
        \State $p_u' \gets p_u$;
    \EndIf
     \If{$p_u'$ and $p_v'$ are in two disconnected $k$-$\widehat{core}$s}
        \State trace two paths starting from $p_u'$ and $p_v'$ up until a common ancestor;
        \State merge pairs of nodes with the same core number on the paths;
    \EndIf
    \EndFunction

\end{algorithmic}}
\end{algorithm}

